\section{Morfismos cartesianos, imágenes inversas, funtores
cartesianos}
\label{VI.5}

Sea~$\cal{F}$ una categoría sobre~$\cal{E}$, con funtor de
proyección~$p$.
\ifthenelse{\boolean{orig}}
{}
{\enlargethispage{.5cm}}

\begin{definition}
\label{VI.5.1} Consideremos un morfismo
$$
\alpha\colon\eta\to\xi
$$
en~$\cal{F}$, y sean
$$
S=p(\xi),\quad T=p(\eta),\quad f=p(\alpha)%
\ifthenelse{\boolean{orig}}{}{.}
$$
Se dice que~$\alpha$ es un \emph{morfismo cartesiano}
\index{cartesiano (morfismo)|hyperpage}%
si, para todo $\eta'\in\Ob(\cal{F}_T)$ y todo $f$-morfismo
$u\colon\eta'\to\xi$, existe un único $T$-morfismo
$\overline{u}\colon\eta'\to\eta$ tal que
$u=\alpha\circ\overline{u}$.
\end{definition}

Esto significa, pues, que para todo $\eta'\in\Ob(\cal{F}_T)$, la
aplicación $v\mto\alpha\circ v$:
\begin{equation}
\tag{i} \Hom_T(\eta',\eta)\to\Hom_f(\eta',\xi)
\end{equation}
es biyectiva. Significa también que el par $(\eta,\alpha)$
\emph{representa el funtor en~$\eta'$} $\cal{F}^{\circ}_T\to\Ens$ del
segundo miembro. Si, para un morfismo dado $f\colon T\to S$ en~$\cal{E}$
y un $\xi\in\Ob(\cal{F}_S)$ dado, existe tal par $(\eta,\alpha)$, es decir,
un morfismo
\marginpar{162}
cartesiano~$\alpha$ en~$\cal{F}$ de destino~$\xi$, tal que
$p(\alpha)=f$, entonces~$\eta$ queda determinado en~$\cal{F}_T$ salvo
isomorfismo único. Se dice entonces que \emph{existe la imagen inversa
de~$\xi$ por~$f$}, y un objeto~$\eta$ de~$\cal{F}_T$ provisto de un
$f$-morfismo cartesiano $\alpha\colon\eta\to\xi$ se llama \emph{una imagen
inversa de~$\xi$ por~$f$}.
\index{imagen inversa|hyperpage}%
A menudo se supone
\ifthenelse{\boolean{orig}}{elegido}{elegida}
tal imagen inversa cada vez que existe (con~$\cal{F}$ fijada); se denotará
entonces la imagen inversa mediante símbolos como~$f^*_{\cal{F}}(\xi)$, o
simplemente~$f^*(\xi)$
\label{indnot:fo}\oldindexnot{$f^*_{\cal{F}}(\xi)$ o $f^*(\xi)$|hyperpage}%
\ifthenelse{\boolean{orig}}{donde~$\xi\times_S T$}
{o bien~$\xi\times_S T$}
cuando estas notaciones no ocasionen
\ifthenelse{\boolean{orig}}{unas confusiones;}{confusión;}
el morfismo canónico $\alpha\colon\eta\to\xi$ se denotará, en lo que
sigue, por~$\alpha_f(\xi)$.
\label{indnot:fp}\oldindexnot{$\alpha_f(\xi)$|hyperpage}%
Si para todo $\xi\in\Ob(\cal{F}_S)$ existe la imagen inversa de~$\xi$
por~$f$, se dirá también que \emph{existe el funtor de imagen inversa
por~$f$ en~$\cal{F}$}, y~$f^*(\xi)$ se convierte entonces en un
\emph{funtor covariante en~$\xi$}, de~$\cal{F}_S$ en~$\cal{F}_T$. Esto
procede de que el segundo miembro de~(i) depende covariantemente de~$\xi$,
es decir, más precisamente, designa un funtor de
$\cal{F}^{\circ}_T\times\cal{F}_S$ en~$\Ens$. Esta dependencia funtorial
de~$f^*(\xi)$ se explicita así: consideremos $f$-morfismos cartesianos
$$
\alpha\colon\eta\to\xi,\quad\alpha'\colon\eta'\to\xi'
$$
y un $S$-morfismo $\lambda\colon\xi\to\xi'$; entonces existe un único
$T$-morfismo $\mu\colon\eta\to\eta'$ tal que
$$
\alpha'\mu=\lambda\alpha
$$
(como se deduce de que~$\alpha'$ es cartesiano).

Señalemos también el hecho inmediato siguiente: consideremos un diagrama
conmutativo
$$
\xymatrix@C=1.2cm{
\xi\ar[d]_-{\lambda}&\ar[l]_-{\alpha}\eta\ar[d]^-{\mu}\\
\xi'&\ar[l]_-{\alpha'}\eta'
}
$$
en~$\cal{F}$,
\marginpar{163}
donde~$\alpha$ y~$\alpha'$ son $f$-morfismos,~$\lambda$ un
$S$-isomorfismo y~$\mu$ un $T$-isomorfismo. \emph{Para que~$\alpha$ sea
cartesiano, es necesario y suficiente que~$\alpha'$ lo sea}.

\begin{definition}
\label{VI.5.2} Un $\cal{E}$-funtor $F\colon\cal{F}\to\cal{G}$ se llama
\emph{funtor cartesiano}
\index{cartesiano (funtor)|hyperpage}%
si transforma los morfismos cartesianos en morfismos cartesianos.
Se denota por $\SheafHom_\cart(\cal{F},\cal{G})$
\label{indnot:fq}\oldindexnot{$\SheafHom_\cart(\cal{F},\cal{G})$|hyperpage}%
la subcategoría plena de
$\SheafHom_{\cal{E}/-}(\cal{F},\cal{G})$ formada por los funtores
cartesianos.

Por ejemplo, considerando~$\cal{E}$ como una categoría sobre~$\cal{E}$
mediante el funtor identidad, todo morfismo de~$\cal{E}$ es cartesiano; por
tanto, un funtor cartesiano de~$\cal{E}$ en~$\cal{F}$ es un funtor sección
$F\colon\cal{E}\to\cal{F}$ que transforma todo morfismo de~$\cal{E}$ en
un morfismo cartesiano; tal funtor se llama una \emph{sección cartesiana}
de~$\cal{F}$ sobre~$\cal{E}$.
\end{definition}

\begin{proposition}
\label{VI.5.3} \textup{(i)}~Un funtor $F\colon\cal{F}\to\cal{G}$ que sea
una $\cal{E}$-equivalencia es un funtor cartesiano.
\textup{(ii)}~Sean~$F,G$ dos $\cal{E}$-funtores
\emph{isomorfos}~$\cal{F}\to\cal{G}$. Si uno es cartesiano, el otro también
lo es. \textup{(iii)}~El compuesto de dos funtores cartesianos
$\cal{F}\to\cal{G}$ y $\cal{G}\to\cal{H}$ es un funtor cartesiano.
\end{proposition}

La afirmación~(iii) es trivial por la definición; (ii) se deduce de la
observación que precede a~\Ref{VI.5.2}; (i) se deduce fácilmente de la
definición y del criterio~\Ref{VI.4.2}~(iii). Más precisamente, un
morfismo~$\alpha$ en~$\cal{F}$ es cartesiano si y solo si~$F(\alpha)$ lo
es.

\begin{corollaire}
\label{VI.5.4}
Sea $F\colon\cal{F}\to\cal{G}$ una $\cal{E}$-equivalencia. Entonces, para
toda categoría~$\cal{H}$ sobre~$\cal{E}$, los funtores correspondientes
$G\mto G\circ F$ y $G\mto F\circ G$ inducen equivalencias de categorías:
\begin{eqnarray*}
\SheafHom_\cart(\cal{G},\cal{H})\lto{\approx}
\SheafHom_\cart(\cal{F},\cal{H})\\
\SheafHom_\cart(\cal{H},\cal{F})\lto{\approx}
\SheafHom_\cart(\cal{H},\cal{G})
\end{eqnarray*}
\end{corollaire}

Esto se deduce, de la manera habitual, del criterio~(i) de~\Ref{VI.4.2} y
de~\Ref{VI.5.3}~(i), (ii) y~(iii).
\marginpar{164}
Puede precisarse que \emph{el $\cal{E}$-funtor
$G\colon\cal{G}\to\cal{H}$ es cartesiano si y solo si $G\circ F$ lo es},
y, análogamente, \emph{un $\cal{E}$-funtor
$G\colon\cal{H}\to\cal{F}$ es cartesiano si y solo si $F\circ G$ lo es}.

De
\ifthenelse{\boolean{orig}}{\Ref{VI.5.4}~(iii)}{\Ref{VI.5.3}~(iii)}
se deduce que, si se considera la subcategoría~$\Cat^\cart_{/\cal{E}}$
\label{indnot:fr}\oldindexnot{$\Cat^\cart_{/\cal{E}}$|hyperpage}%
de~$\Cat_{/\cal{E}}$ cuyos objetos son los mismos que los de
$\Cat_{/\cal{E}}$ y cuyos morfismos son los funtores \emph{cartesianos},
se tienen, como en el~\No \Ref{VI.2}, operaciones de composición:
$$
\SheafHom_\cart(\cal{F},\cal{G})\times
\SheafHom_\cart(\cal{G},\cal{H})\to\SheafHom_\cart(\cal{F},\cal{H})
$$
\ifthenelse{\boolean{orig}}{inducida}{inducidas}
por las del~\No \Ref{VI.2}, que permiten considerar
$\SheafHom_\cart(\cal{F},\cal{G})$ como un funtor en~$\cal{F}$,
$\cal{G}$, de la categoría $\left(\Cat^\cart_{/\cal{E}}\right)^\circ
\times\Cat^\cart_{/\cal{E}}$
\ifthenelse{\boolean{orig}}{en~$\Cat$).}{en~$\Cat$.}
Necesitaremos esta observación sobre todo para el caso en que
$\cal{F}=\cal{G}$:

\begin{definition}
\label{VI.5.5}
Sea $\cal{F}$ una categoría sobre~$\cal{E}$. Se denota por
$$
\varprojLim\cal{F}/\cal{E}
$$
\label{indnot:fs}\oldindexnot{$\varprojLim\cal{F}/\cal{E}$|hyperpage}%
la categoría de los $\cal{E}$-funtores cartesianos
$\cal{E}\to\cal{F}$, es decir, de las secciones cartesianas de~$\cal{F}$
sobre~$\cal{E}$.
\end{definition}

De acuerdo con lo que acaba de decirse,
$\varprojLim\cal{F}/\cal{E}$ es un funtor en~$\cal{F}$, de la categoría
$\Cat^\cart_{/\cal{E}}$ en la categoría~$\Cat$.

Veremos más adelante las relaciones entre esta operación~$\varprojLim$ y la
noción de límite proyectivo de categorías, así como numerosos ejemplos.
